% page 322 du fac-simile (image p-321 du scan)
% bloc 152.4 x 233.3 mm ; marge gauche 8.0 mm
\pgc{-12.700mm}{0.000mm}{
\l{\cel{5}314}
\l{}
\l{\sou{kresto.}-\cel{2}I. Surkreskajo karnoza, di qua la konturo esas}
\l{\cel{3}sesgura dento-forme, qua salias sur la kapo di ula}
\l{\cel{3}galinacei. (\sou{analogie}: sur la kapo di ula repteri, fi-}
\l{\cel{3}shi, uceli.)\cel{2}- II. Parto salianta, streta, sur la so-}
\l{\cel{3}mito di monto, di ondo, la supra parto di tekto, di}
\l{\cel{3}muro. - EFIS.}
\l{}
\l{\sou{krestomatio}. Selektajo ek texto-peci extraktita de la verki}
\l{\cel{3}da la autori klasika. - DEFIRS.}
\l{}
\l{\sou{kreto}. Kalko-karbonato amorfa, quan onu renkontras forme di}
\l{\cel{3}strati plu o min dika, en ula sulo-strati. -\cel{1}\sou{Kreto-kra-}}
\l{\cel{3}\sou{yono}\cel{1}: bastoneto ek ta substanco, per qua onu skribas}
\l{\cel{3}sur tabelo nigra o obskure-verda (en skolo). - DEFIS.}
\l{}
\l{kretino.\cel{2}(\sou{patol.)}\cel{2}I. Persono rakitika, idioto maxim-multa-}
\l{\cel{3}kaze strumoza. - II. (\sou{metaf.}) Persono stupida quale kre-}
\l{\cel{3}tino. - DEFIRS.}
\l{}
\l{\sou{kretono.}\cel{2}Telo rezistiva, kun warpo ek kanabo,e wefto ek}
\l{\cel{3}lino. - DEFRS.}
\l{}
\l{\sou{krevar}.\cel{2}I. (\sou{netrans}.)Apertar su, spliteskante, pro tens-}
\l{\cel{3}eso ecesanta. - II. (\sou{trans.)}\cel{3}Apertar, splitante. - FI}
\l{}
\l{\sou{kreveto}. (\sou{zool.)}\cel{2}Mikra krustaceo dekapeda, manjebla, quan}
\l{\cel{3}onu renkontras sur la litoro marala. - L.\cel{1}\sou{crangon vul-}}
\l{\cel{3}\sou{garis}. - F.}
\l{}
\l{\sou{krevisar}.\cel{2}I. (\sou{netrans.})\cel{2}Subisar fendeso plu o min profun-}
\l{\cel{3}da (\sou{aludante la surfaco di korpo}).- II. (\sou{trans.}) Fendar}
\l{\cel{3}plu o min profunde (la surfaco di korpo). - EFI.}
\l{}
\l{\sou{kriar}. (\sou{net}rans). Lansar voco-sonuni stridanta e neartiku-}
\l{\cel{3}lita.- EFIRS.}
\l{}
\l{\sou{kriblo}. Recevuyo di qua la kulo esas borita ye trui inter-}
\l{\cel{3}egala, por separer la objekti di groseso neinteregala,}
\l{\cel{3}lasanta trapasar ici, retenante iti. - FIS.}
\l{}
\l{\sou{krik}\cel{1}!\cel{2}Interjeciono qua imitas la bruiso de kozo lacerata.}
\l{}
\l{\sou{kriko}. Mashino por sublevar nealte kargajo (quadri, e c.)}
\l{\cel{3}per dento-stango quan elevas roto dentoza, movata per}
\l{\cel{3}manivelo. - FIS.}
\l{}
\l{\sou{krik-krak}\cel{1}!\cel{3}Interjeciono per qua onu probas reproduktar}
\l{\cel{3}la bruiso sika di kozo qua ruptesas, laceresas.}
\l{}
\l{\sou{kriketo.}\cel{2}Ludilo, origine Angla, e qua la ludanti, dividi-}
\l{\cel{3}ta a du partisi, lansas singlu-suafoye baloneti grava,}
\l{\cel{3}per kroso ligna. - EF.}
}
